% Build the layer, Chapter 4, first build: a line of a digest.
% Inserted at the end of \S\ref{sec:zoom}, after the scene "Alice, 3:32
% p.m." The protocol there says how a line of a digest is to be made. This
% makes one in the kernel's decision core, and makes landing wait for it.
%
% No code or output in this file is typed. Each listing is a file under
% excerpts/, which the kernel's tools/excerpts.mjs cuts from its source and
% its recorded runs (see figures/pd-listing.tex). The figure draws one of those recorded tables as
% it was printed. To change a listing or the figure, change the kernel and
% cut again. The prose names functions, never line numbers, so that a
% listing which moves in its file does not make a sentence wrong.
%
% \pdlistingroom{n} before a listing of n lines keeps it on one page. When a
% listing changes length, change the number with it.
\subsection{Build the layer: a line that leads back to the work}
\label{sec:ls-line}

\noindent\emph{Express lane: the line is built three times. The first
prints what the worker says, and Alice lands a change that failed the test
it was to keep. The second asks for links and waits until each is opened.
One worker attaches the log of other work, and another changes its work
after Alice has read it. The third is made by the decision core from the
kernel's record, and landing makes it again and waits for the hashes on
it. A walk of every afternoon of one piece of work counts 10,094 states in
which the second surface lets work land, and 1,214 of them are right. Under
the kernel's surface all 1,092 are.}

The protocol above says how a line is to be made, and the scene after it
has Alice follow it by eye. \pdkernelname{} builds a surface that follows
it, for one piece of work and one act that cannot be undone.

The work is the checkout repair of \S\ref{sec:mara}, which starts at
revision~7. A worker changes it, and each change makes the next revision.
The kernel stores every diff. It also runs the test that the job says to
keep, which this section calls the \emph{kept test}, and a change that
deletes that test fails it. The worker may say things about its work. Alice
reads one line about the work on a read-surface and opens what the line
points to. Then she lands the work or leaves it.

Every run below ends on one question. A landing is \emph{right} if a run of
the kept test had passed on the revision that landed, and Alice had been
shown that revision's diff.

\subsubsection{Naive: the line is what the worker says}

The obvious surface asks the author. The worker says how its work went, and
the line is what it said, after its name and the job's
(Listing~\ref{alg:ls-line-said}).

\pdlistingroom{11}
\pdexcerpt[caption={The first try. The line is the worker's words, and
landing waits for nothing.},label={alg:ls-line-said}]{ch04-said}

This is Row~B of the example that opens this section. Nothing stands
behind the line, so \texttt{land} has nothing to wait for.

\subsubsection{Break: all green}

Drake repairs the retry and deletes the kept test
(Listing~\ref{alg:ls-line-run-green}).

\pdlistingroom{14}
\pdrun[caption={A worker that says all green, at the first surface. Nobody
here is asked who they are: the names say who would
act.},label={alg:ls-line-run-green}]{green-said}

The kernel ran the kept test on revision~8, and it failed. That run is in
the kernel's record, and the line does not draw on the record. Alice read a
true account of what Drake said, and she landed work that the record called
broken.

\subsubsection{Break again: a real log of other work}

Invariant~\ref{def:zoom} says what to repair: every summary must reach the
artifact it summarizes. So the second surface refuses a line with nothing
attached. It shows each thing the worker attached as a claim that Alice can
open, and landing waits until she has opened every one
(Listing~\ref{alg:ls-line-linked}).

\pdlistingroom{23}
\pdexcerpt[caption={The second try. A line needs links, and landing waits
for each to be opened.},label={alg:ls-line-linked}]{ch04-linked}

This is Row~A. Drake's bare ``all green'' is refused now, because nothing is
attached to it. So Drake attaches a log
(Listing~\ref{alg:ls-line-run-borrowed}).

The log Alice is shown is real, and the run in it passed. It is a run of
revision~7, and Drake's work stands at revision~8. The example above warned
that a dishonest agent could reuse Row~A's log, and
\S\ref{sec:completionist} says that a passing receipt for another tree
cannot close the obligation. This surface compares nothing. It shows Alice
what the worker chose to attach.

\pdlistingroom{17}
\pdrun[caption={The second break. The link opens, and it leads to a real
log of other work.},label={alg:ls-line-run-borrowed}]{borrowed-linked}

\subsubsection{Break a third time: the work changes after it was read}

The next break needs no false word from anyone
(Listing~\ref{alg:ls-line-run-switch}).

\pdlistingroom{21}
\pdrun[caption={The third break. Alice opened everything on the line, and
the work changed afterwards.},label={alg:ls-line-run-switch}]{switch-linked}

Everything Alice was shown was true of revision~8. Wren then added a
migration, and the kernel's run of revision~9 passed. The surface landed
revision~9. It had kept the numbers of the claims she opened, and claims 1
and 2 were still ticked. She read one revision and landed another. This is
the unrelated migration of \S\ref{sec:mara}, and it arrived after the
review.

\subsubsection{Fix: the line is made from the record}

The breaks share a fault. Someone with a stake in the landing chose what
Alice saw: the words under the first surface, and the links under the
second. The third break adds that the work moved after she had looked.
\pdkernelname{} takes the choice away from the worker. The line is
made by the decision core (\weld{prereq:sec:prereq-build}{decision core}),
which is shown the kernel's record of the work and nothing the worker said.
The record holds the work's revision, the hash of its stored diff, the
kernel's newest run of the kept test and the hashes of what Alice has
opened.

\texttt{render} makes the line from that record. A line is a list of
claims, and a claim is a statement with the SHA-256 of the text that backs
it. The first claim is that the change is this diff. If the kernel has run
this revision, a second claim says that the test passed or failed, with the
hash of the run's log. The line is marked verified only if such a run
passed (Listing~\ref{alg:ls-line-render}). Two guards stand behind it
(Listing~\ref{alg:ls-line-render-guards}).

\pdlistingroom{17}
\pdexcerpt[caption={The line, made from the
record.},label={alg:ls-line-render}]{ch04-render}

\pdlistingroom{18}
\pdexcerpt[caption={The guards on a line. A line marked \texttt{MUTANT} is
compiled only into a build of the core that has that guard removed, and the
verification below runs each of those
builds.},label={alg:ls-line-render-guards}]{ch04-render-guards}

\texttt{backing} makes no line for work with no diff stored.
\texttt{current} lets a run count only if it is a run of the revision the
work stands at. The kernel's own run of an earlier revision is a passing
receipt for another tree, as the log Drake borrowed was, and
\texttt{current} leaves it off the line.

Nothing the worker says is an input to \texttt{render}, which answers the
first break (Listing~\ref{alg:ls-line-run-green-core}).

\pdlistingroom{16}
\pdrun[caption={The first break, on the kernel's
surface.},label={alg:ls-line-run-green-core}]{green-core}

The line says \texttt{NOT VERIFIED}, and it names the failing log by its
hash. The refusal carries a named reason, which is the protocol's fourth
step.

Attaching another worker's log changes nothing, because what a worker
attaches is not in the record
(Listing~\ref{alg:ls-line-run-borrowed-core}).

\pdlistingroom{11}
\pdrunlines[caption={The second break, on the kernel's surface. The claim
Alice opens is the diff of revision~8, and it shows her the
deletion.},label={alg:ls-line-run-borrowed-core}]{borrowed-core}{1,8-17}{1 and 8\textendash 17}

Opening is a decision of the core's as well. The kernel's store hands over
the text it holds under a claim's hash, and \texttt{open} is given that
text. The claim counts as opened only if \texttt{same} agrees. What the
core then records is the claim's hash and not its number
(Listing~\ref{alg:ls-line-open}).

\pdlistingroom{19}
\pdexcerpt[caption={Opening a claim.},label={alg:ls-line-open}]{ch04-open}

\texttt{same} hashes the text again and compares the result with the
claim's own hash (Listing~\ref{alg:ls-line-same}).

\pdlistingroom{9}
\pdexcerpt[caption={The guard on what is shown. What the store hands over
is hashed again.},label={alg:ls-line-same}]{ch04-same}

\texttt{land} answers the third break. It begins with \texttt{read}, which
calls \texttt{render}, so the line it judges is the line as the work stands
at that moment. The work must be verified, and no claim on that line may be
unopened (Listing~\ref{alg:ls-line-land}).

\pdlistingroom{15}
\pdexcerpt[caption={Landing. The line is made again, and every claim on it
must have been opened.},label={alg:ls-line-land}]{ch04-land}

The order of the two checks is the order in which a reader is told things.
Work that is not verified is refused for that, whatever has been opened, so
nobody is sent to read a diff that cannot land. Verified work is refused
with the number of the first claim that is still unopened, which tells
Alice what to open next.

\texttt{unopened} looks for a claim on the line whose hash is not among the
hashes Alice opened, and \texttt{seen} is the comparison it makes
(Listing~\ref{alg:ls-line-land-guards}).

\pdlistingroom{21}
\pdexcerpt[caption={The guards on landing.},label={alg:ls-line-land-guards}]{ch04-land-guards}

When Wren changes the work after the review, the new line names a new diff,
and Alice has not opened it (Listing~\ref{alg:ls-line-run-switch-core}).

\pdlistingroom{22}
\pdrun[caption={The third break, on the kernel's
surface.},label={alg:ls-line-run-switch-core}]{switch-core}

The same guard answers the unread pile of \S\ref{sec:mara}. Verified work
that nobody has opened does not land, and it lands once it has been opened
(Listing~\ref{alg:ls-line-run-pile}).

\pdlistingroom{13}
\pdrunlines[caption={Verified work, landed after it was opened and not
before. The steps left out are the first four of
Listing~\ref{alg:ls-line-run-switch-core}.},label={alg:ls-line-run-pile}]{pile-core}{1,11-22}{1 and 11\textendash 22}

This is the forced zoom of the protocol's third step. The build has no flag
for it. Its one act cannot be undone, so every line would carry the flag
set.

\keyidea[The reader's line]{Whoever chooses what a summary points to
decides what its reader sees. So the line is made from the kernel's record,
by code with no stake in the landing, and it names each backing text by its
hash. Work can change between a reading and an act, so the act makes the
line again and waits for the hashes that are on it then.}

\Needspace{12\baselineskip}
\subsubsection{Verify: every afternoon of one piece of work}

\paragraph{Every afternoon.}
The runs above are seven afternoons. The kernel's tests walk every
afternoon of one piece of work within small bounds. The worker changes the
work up to 2 times, well or badly, and says up to 2 things of it. Its words
come bare or with something attached: a log of other work, or its own diff
and newest log. The kernel runs the kept test whenever it gets to it. The
store goes wrong at most once, which means that the text it holds under the
diff's hash becomes some other text. Alice opens up to 3 claims, and she
may try to land in any state.

Every state is held to rules written apart from the core, which do their
own hashing. A line \emph{says more than it shows} if it tells its reader
that the work is good and carries no claim that leads to a passing run of
the work as it stands. A landing is counted as not verified if no run of
its revision had passed, and as not read if no text Alice was shown hashes
to its diff (Listing~\ref{alg:ls-line-walk}).

\pdlistingroom{22}
\pdrun[caption={Every afternoon of one piece of work, under each
surface.},label={alg:ls-line-walk}]{lines}

The first surface shows a line in each of its 2,388 states, and 2,160 of
those lines say more than they show. It lets the work land in every state.
The second surface has 10,094 states in which work can land. In 7,002 of
them no run of the work had passed, and in 6,544 Alice had not been shown
its diff. The kernel's surface shows 14,308 lines, and none says more than
it shows. It lets work land in 1,092 states, each of them verified and
read, and it refused no landing that was both.

\paragraph{Every landing.}
The two faults overlap, so the last two columns of the listing cannot be
added. Figure~\ref{fig:ls-line-landings} sorts every landing into one of
four kinds.

% Brief (grammar: small multiples, one bar for each surface, cut once
%   across and once down by the two halves of the landing rule; atlas
%   I/fig:ls-line-landings).
% Reader question: of the landings each surface allows, how many are right,
% and what is wrong with the others?
% Claim: none of the first surface's 2,388, then 1,214 of the second's
% 10,094, and all 1,092 of the kernel's. The wrong ones are of two kinds
% that overlap: work no run had passed, and work whose diff Alice had not
% been shown.
% Evidence: every landing in every afternoon of one piece of work, within
% the bounds the walk prints [internal, harbor lines --each; session
% lines-each.txt, drawn as printed, and
% tests/harbor-research/test_ls_line_numbers.py holds the counts in the
% caption and the prose to it].
% Must distinguish: a landing that was right from one that was not; work
% that was not verified from work that was not read; one surface from the
% next.
% Grammar chosen: each bar is every landing one surface allowed. It is cut
% across by whether a run of the work had passed, and each part is cut
% down by whether Alice had been shown the diff, so that an area is a share
% of that surface's landings and the four kinds keep their places from bar
% to bar.
% Grammar rejected: one stacked bar for each surface, which sets the four
% kinds in a row and hides that two questions make them; a bar on one
% scale for all landings, which compares counts of states from walks of
% different sizes.
% Five-second test: which bar is black from end to end?
\begin{figure}[H]
\centering
\pdfigurehue{pdcobalt}%
\begin{tikzpicture}[pd figure,x=1cm,y=1cm,
    /pgf/number format/.cd,int detect,1000 sep={,},assume math mode=true,
    /tikz/.cd,
    kind/.style={pd label,text=white,inner sep=0pt}]
  % The table is a recorded run, read as it was printed: the first line is
  % the command, and every remark starts with '#'.
  \pgfplotstableread[skip first n=1]{excerpts/lines-each.txt}\pdlandings
  \def\pdbarwide{10.2}\def\pdbarhigh{1.6}\def\pdbarstep{2.45}
  % One cell of a bar: #1 and #2 are its corners, #3 its style, #4 its name
  % and #5 its count. A kind that no landing was of is not drawn.
  \newcommand{\pdlandingcell}[5]{%
    \ifnum#5>0
      \fill[#3] #1 rectangle #2;
      \draw[draw=white,line width=1.1pt] #1 rectangle #2;
      \path #1 -- node[kind] {#4 \pgfmathprintnumber{#5}} #2;
    \fi}
  \foreach \row/\name in {0/{the first surface},1/{the second surface},2/{the kernel's surface}}{
    \pgfplotstablegetelem{\row}{landings}\of\pdlandings \let\pdall\pgfplotsretval
    \pgfplotstablegetelem{\row}{right}\of\pdlandings \let\pdfine\pgfplotsretval
    \pgfplotstablegetelem{\row}{unread}\of\pdlandings \let\pdunread\pgfplotsretval
    \pgfplotstablegetelem{\row}{unverified}\of\pdlandings \let\pdunverified\pgfplotsretval
    \pgfplotstablegetelem{\row}{neither}\of\pdlandings \let\pdneither\pgfplotsretval
    \pgfmathsetmacro{\pdtop}{-\row*\pdbarstep}
    \pgfmathsetmacro{\pdfoot}{\pdtop-\pdbarhigh}
    % Across: the landings of work that a run had passed, then the others.
    \pgfmathtruncatemacro{\pdpassed}{\pdfine+\pdunread}
    \pgfmathtruncatemacro{\pdfailed}{\pdunverified+\pdneither}
    \pgfmathsetmacro{\pdcut}{\pdpassed/\pdall*\pdbarwide}
    % Down each part: the landings she had been shown the diff of, then
    % the others.
    \ifnum\pdpassed>0
      \pgfmathsetmacro{\pdmid}{\pdtop-\pdfine/\pdpassed*\pdbarhigh}
      \pdlandingcell{(0,\pdtop)}{(\pdcut,\pdmid)}{fill=pdink}{right}{\pdfine}
      \pdlandingcell{(0,\pdmid)}{(\pdcut,\pdfoot)}{fill=pderror}{not read}{\pdunread}
    \fi
    \ifnum\pdfailed>0
      \pgfmathsetmacro{\pdmid}{\pdtop-\pdunverified/\pdfailed*\pdbarhigh}
      \pdlandingcell{(\pdcut,\pdtop)}{(\pdbarwide,\pdmid)}{fill=pderror}{not verified}{\pdunverified}
      \pdlandingcell{(\pdcut,\pdmid)}{(\pdbarwide,\pdfoot)}{fill=pderror}{neither}{\pdneither}
    \fi
    % A line inside as well as a colour, so that right differs from wrong
    % in grey.
    \ifnum\pdfine>0
      \pgfmathsetmacro{\pdmid}{\pdtop-\pdfine/\pdpassed*\pdbarhigh}
      \draw[draw=white,line width=.5pt] (.09,\pdtop-.09) rectangle (\pdcut-.09,\pdmid+.09);
    \fi
    \node[pd title,anchor=south west,inner sep=0pt] at (0,\pdtop+.13)
      {\name: \pgfmathprintnumber{\pdfine} of \pgfmathprintnumber{\pdall} landings were right};
  }
  \node[pd label,anchor=north,align=center] at (\pdbarwide/2,-2*\pdbarstep-\pdbarhigh-.25)
    {across: a run of the work had passed, then none had\\
     down: Alice had been shown its diff, then she had not};
\end{tikzpicture}
\caption{\textbf{Of 10,094 landings under the surface with links, 1,214
were right.} Each bar is every landing one surface allowed, and an area is
a share of them. Black is a landing of verified work whose diff Alice had
been shown.
[internal, \texttt{harbor lines -{}-each}]}
\label{fig:ls-line-landings}
\end{figure}


Under the first surface no landing is right, because it has nothing to
open. In 648 of them the work happened to be verified, and in the other
1,740 it was not. Under the second surface 1,214 of 10,094 landings are
right, which is about one in eight. In 4,666 the work was neither verified
nor read. Each of those landings had links on its line, and the surface
counted every one as opened. Under the kernel's surface the bar has one
kind.

\paragraph{Any record.}
The walk stops at 2 changes and 3 openings. Kani lifts the bounds on the
record (\weld{anchor:sec:correct}{Kani}). It compiles the core's own
\texttt{land} and asks a solver about every record at once
(Listing~\ref{alg:ls-line-kani}).

\pdlistingroom{13}
\pdexcerpt[caption={A Kani harness on landing. The record is any record,
with any two hashes opened.},label={alg:ls-line-kani}]{ch04-kani}

\texttt{any\_work} builds any record: any revision, a diff or none, a run
of any revision or none, and any two hashes opened. If \texttt{land} admits
it, the harness's 2 assertions say that the line is verified and that every
claim on it is among the hashes opened. Its \texttt{cover} line makes the
solver find a record that does land. Without it the assertions could hold
for want of any case to try.

The harness pays for its reach. Under Kani a hash is one byte and not a
string of 64 characters. With strings, the solver ran out of 8~GB of memory
on a harness of a few lines. The rules compare hashes and never look inside
one, so they decide the same way on bytes. The code Kani checked therefore
differs from the code that runs by one type, and SHA-256 is left out of it
(Listing~\ref{alg:ls-line-kani-run}).

\pdlistingroom{9}
\pdrunlines[caption={Kani~0.68 on the core's rules for a
line.},label={alg:ls-line-kani-run}]{kani-zoom}{1-9}{1\textendash 9}

The first verdict is a control. It puts the question to a gate that counts
how many things were opened and compares no hashes, and that harness fails
as it must.

\paragraph{Each guard removed.}
This build has six guards. Five are in the listings above, and the sixth
is the \texttt{verified} that \texttt{land} calls. The kernel has a build
with each one taken out. The harnesses are run on each
(Listing~\ref{alg:ls-line-kani-guards}), and so is the walk
(Listing~\ref{alg:ls-line-guards-walk}).

\pdlistingroom{14}
\pdrunlines[caption={The harnesses again, on six builds of the core with
one guard removed from each.},label={alg:ls-line-kani-guards}]{kani-zoom}{1,10-22}{1 and 10\textendash 22}

Read the walk's rows against the breaks. A core that counts any run is the second
break inside the core: 1,716 of its lines say more than they show, and 462
of its landings are of work no run had passed. A core that counts what was
opened is the third break, with 912 landings of unread work. A core that
makes a line with no diff behind it lands nothing wrong, and 20 of its
lines lead to nothing in the store.

\pdlistingroom{26}
\pdrun[caption={The walk again, on the same six
builds.},label={alg:ls-line-guards-walk}]{lines-guards}

The row to study is \texttt{trust-store}. A core that does not hash what
the store hands over lands unread work in 376 states of the walk, and no
harness fails on it. That is the price paid above. Under Kani the
comparison of a text with its hash is the solver's to answer, so the guard
is not in the code Kani sees. The gate that
\pdchapref{sealed}{The Sealed Harbor} builds
(\weld{sealed:sec:sealed-gate}{walk of pairs}) had a guard that its walk
could not see and Kani could. Here it is the other way round. Each tool is blind somewhere, and running both on every build with a
guard removed is how the blind places are found.

\pdscope{The walk covers one piece of work, within the bounds that
Listing~\ref{alg:ls-line-walk} prints: 2 changes, 2 sayings, 3 openings and
one fault of the store. Its counts are of states, with no
chances attached: a careless afternoon and a careful one each count once.
The harnesses cover one decision of the core's on any record, with a byte
standing for a hash.}

\begin{pdboundary}
\setlength{\parskip}{4pt plus 1pt}
A true line is not a right change. The line says that the kept test passed
on this revision, and landing says that Alice was shown this diff. Wren's
revision~9 passes the kept test and carries a migration nobody asked for.
If Alice opens its diff she is shown the migration, and whether she sees it
is nothing the core can know. The job has one kept test, and nothing here
asks whether it is the right test.

Opened is not read. The core records that a text was shown and that it
hashed to the claim. An operator who opens every claim without looking
lands what a careful one would not. How much reading an operator can do is
the subject of \S\ref{sec:operator}, and of this chapter's next build.

Nobody is asked who they are. A move names no one, so a worker that can
reach the core can open its own claims. The kernel's run of the kept test
is taken as given, and a runner that a worker controls is outside the
model.

The store has to keep every text a line has named, for as long as the line
may be read. If it changes a text before Alice has opened it, the claim
cannot be opened. The work behind it then does not land, however good it
is.
\end{pdboundary}
